%% appA --- Sciagawka C# do Pythona
%%
%% Napisany we wrzesniu 2026. Jedna tabela na czesc ksiazki. Kazdy wiersz
%% nazywa rozdzial, ktory go UCZY, przez \ref zamiast wpisanego numeru, a
%% tools/check_structure.py --cheatsheet oblewa, gdy wiersz wskazuje na
%% etykiete, ktorej nie ma, albo na rozdzial, ktory jest wciaz zaslepka.
%% Wiersz, ktorego nic nie uczy, jest wada tego dodatku, a nie faktem o
%% ksiazce.

\chapter{Ściągawka C\# do Pythona}
\label{app:cheatsheet}

\noindent
Wszystko tutaj jest gdzieś w książce nauczane, a ostatnia kolumna mówi
gdzie. To właśnie jest sens tej kolumny: wiersz jest twierdzeniem o
rozdziale, a wiersz, którego żaden rozdział nie uczy, jest wadą tego
dodatku, a nie luką w twojej wiedzy.

Czytaj dwie środkowe kolumny jako tłumaczenie, a nie jako równoważność.
Kilka wierszy odwzorowuje się czysto; kilka odwzorowuje się z różnicą,
która coś kosztuje, i to na nie rozdział poświęca swoje strony. Tam, gdzie
wiersz sam z siebie wprowadzałby w błąd, mówi o tym w kolumnie Pythona.

\section{Część I \dash{} Środowisko uruchomieniowe i narzędzia}
\label{sec:appA-part1}

\begin{cheatsheet}{C\#}{Python}{Rozdz.}
CLR, JIT do kodu natywnego & CPython, kod bajtowy w
  \code{\_\_pycache\_\_}; brak JIT-a, na którym można polegać &
  \ref{ch:runtime} \\
IL, \code{ildasm} & kod bajtowy, \api{dis.dis} na funkcji &
  \ref{ch:runtime} \\
wątki skalują pracę procesora & nie skalują: blokada serializuje kod
  bajtowy. Skalują pracę I/O & \ref{ch:runtime} \\
\code{static void Main} & \code{if \_\_name\_\_ == "\_\_main\_\_"} &
  \ref{ch:runtime} \\
\code{Parallel.For} & \api{multiprocessing} albo budowa wolnowątkowa &
  \ref{ch:runtime} \\
NuGet & PyPI, a do rozmowy z nim \pkg{uv} & \ref{ch:packaging} \\
\code{packages.lock.json} & \code{uv.lock}, commitowany &
  \ref{ch:packaging} \\
\code{dotnet restore} & \code{uv sync --locked} & \ref{ch:packaging} \\
\code{.csproj} & \code{pyproject.toml} & \ref{ch:packaging} \\
\code{dotnet run} & \code{uv run}, które rozwiązuje środowisko na każde
  wywołanie & \ref{ch:packaging} \\
narzędzie globalne & środowisko na projekt; nie ma dobrej instalacji
  globalnej & \ref{ch:packaging} \\
kompilator sprawdza typy & checker sprawdza, jeśli go zaprosisz: pyright
  albo mypy & \ref{ch:typing} \\
\code{int?}, \code{string?} & \code{int | None}, \code{str | None} &
  \ref{ch:typing} \\
\code{List<int>}, \code{Dictionary<K,V>} & \code{list[int]},
  \code{dict[K, V]} & \ref{ch:typing} \\
interfejs & \code{Protocol}, spełniany strukturalnie, bez żadnej
  deklaracji & \ref{ch:typing} \\
\code{where T : IComparable} & ograniczenie na parametrze typu, składnia
  PEP~695 & \ref{ch:typing} \\
\code{dynamic} & \code{Any} \dash{} ale propaguje się i ucisza checker
  dalej w łańcuchu & \ref{ch:typing} \\
wzorzec \code{is T t} & \api{isinstance}, a \code{TypeIs}, żeby powiedzieć
  to checkerowi & \ref{ch:typing} \\
\end{cheatsheet}

\section{Część II \dash{} Język, odwzorowany}
\label{sec:appA-part2}

\begin{cheatsheet}{C\#}{Python}{Rozdz.}
\code{Equals} + \code{GetHashCode} & \code{\_\_eq\_\_} +
  \code{\_\_hash\_\_}, a pominięcie drugiej usuwa hashowanie &
  \ref{ch:objects} \\
\code{ToString()} & \code{\_\_repr\_\_}, a opcjonalnie
  \code{\_\_str\_\_} & \ref{ch:objects} \\
przeciążanie operatorów & \code{\_\_add\_\_}, \code{\_\_lt\_\_} i reszta,
  żadna nie zadeklarowana & \ref{ch:objects} \\
\code{IComparable} & \code{\_\_lt\_\_} albo \code{order=True} na
  dataclassie & \ref{ch:objects} \\
\code{record}, \code{with} & \code{@dataclass(frozen=True)},
  \api{dataclasses.replace} & \ref{ch:objects} \\
właściwość & \code{@property}, awansowane ze zwykłego atrybutu, gdy jest
  potrzebne & \ref{ch:objects} \\
indekser, \code{Count} & \code{\_\_getitem\_\_}, \code{\_\_len\_\_} &
  \ref{ch:objects} \\
\code{IEnumerable<T>} & \code{\_\_iter\_\_} & \ref{ch:objects} \\
\code{enum}, \code{[Flags]} & \code{Enum}, \code{StrEnum}, \code{Flag} &
  \ref{ch:objects} \\
pole \code{static} & atrybut klasy \dash{} ale czytelny przez
  \code{self}, i to jest pułapka & \ref{ch:objects} \\
\code{struct} & \code{\_\_slots\_\_} jest najbliższą rzeczą i nie jest
  typem wartościowym & \ref{ch:objects} \\
\code{/} na dwóch \code{int}-ach & \code{//}, które zaokrągla w dół, a nie
  obcina & \ref{ch:objects} \\
\code{decimal} & \api{decimal.Decimal} & \ref{ch:objects} \\
\code{Func<>}, delegat & funkcja z adnotacją
  \code{Callable[[A], B]} & \ref{ch:functions} \\
atrybut, który czyta framework & nic: dekorator \emph{działa} &
  \ref{ch:functions} \\
middleware wokół wywołania & dekorator & \ref{ch:functions} \\
\code{params}, opcjonalne, nazwane & \code{*args}, wartości domyślne,
  słowa kluczowe i znaczniki \code{/} oraz \code{*} & \ref{ch:functions} \\
LINQ, odroczone i powtarzalne & wyrażenie listowe (zachłanne) albo
  generatorowe (odroczone, jednoprzebiegowe) & \ref{ch:functions} \\
\code{yield return} & \code{yield}, plus \code{send} i wartość zwracana &
  \ref{ch:functions} \\
\code{using} / \code{IDisposable} & \code{with} /
  \api{contextlib.contextmanager} & \ref{ch:functions} \\
wyrażenie \code{switch} & \code{match}, ze wzorcami sekwencji i mapowania &
  \ref{ch:functions} \\
\code{try} / \code{catch} / \code{finally} & \code{try} / \code{except} /
  \code{finally} / \code{else} & \ref{ch:errors} \\
\code{catch \{ \}} & najwyżej \code{except Exception:}; gołe
  \code{except:} łapie też przerwanie & \ref{ch:errors} \\
\code{throw;} & gołe \code{raise} & \ref{ch:errors} \\
\code{throw new X(..., inner)} & \code{raise X(...) from err} &
  \ref{ch:errors} \\
\code{AggregateException} & \code{ExceptionGroup} i \code{except*} &
  \ref{ch:errors} \\
\code{FooException} & \code{FooError} \dash{} N818 w ruffie &
  \ref{ch:errors} \\
\code{TryGetValue} & proś o wybaczenie: \code{try} /
  \code{except KeyError} & \ref{ch:errors} \\
\code{using X;} & \code{import x} albo \code{from x import name}, nigdy
  \code{import *} & \ref{ch:imports} \\
zestaw (assembly) & pakiet; moduł to plik, który wykonuje się raz &
  \ref{ch:imports} \\
\code{internal} & wiodące podkreślenie, wyłącznie umownie &
  \ref{ch:imports} \\
\code{IServiceCollection} & korzeń kompozycji: funkcja,
  \api{functools.partial} i \code{Protocol} & \ref{ch:imports} \\
\code{Program.cs} & punkt wejścia \code{[project.scripts]}, który pisze
  \code{uv init} & \ref{ch:imports} \\
\end{cheatsheet}

\section{Część III \dash{} Współbieżność}
\label{sec:appA-part3}

\begin{cheatsheet}{C\#}{Python}{Rozdz.}
\code{Task} & korutyna \dash{} ale zimna: wywołanie jej nic nie
  uruchamia & \ref{ch:asyncio} \\
\code{Task.Run} & \api{asyncio.to\_thread} & \ref{ch:asyncio} \\
\code{Task.WhenAll} & \api{asyncio.gather} albo \code{TaskGroup}, jeśli
  chcesz anulowanego rodzeństwa & \ref{ch:asyncio} \\
\code{CancellationToken} & \code{CancelledError}, dostarczany przy
  \code{await}; nic nie jest odpytywane & \ref{ch:asyncio} \\
\code{CancelAfter} & \api{asyncio.timeout} & \ref{ch:asyncio} \\
\code{ConfigureAwait(false)} & nic: nie ma kontekstu synchronizacji &
  \ref{ch:asyncio} \\
\code{.Result}, \code{.Wait()} & nic bezpiecznego: blokujące wywołanie
  zamraża całą pętlę & \ref{ch:asyncio} \\
\code{SemaphoreSlim} & \api{asyncio.Semaphore} & \ref{ch:asyncio} \\
\code{IAsyncEnumerable<T>} & \code{async for} po generatorze
  asynchronicznym & \ref{ch:asyncio} \\
\code{HttpClient} & \code{httpx.AsyncClient} \dash{} jedna instancja, i
  ma już termin & \ref{ch:asyncio} \\
pula wątków rośnie & domyślny executor ma sufit, i to niski &
  \ref{ch:asyncio} \\
\end{cheatsheet}

\section{Część IV \dash{} Wdrażanie}
\label{sec:appA-part4}

\begin{cheatsheet}{C\#}{Python}{Rozdz.}
ASP.NET Core & FastAPI na uvicornie & \ref{ch:services} \\
kontroler & router i funkcja z dekoratorem & \ref{ch:services} \\
wiązanie modelu & model pydantica w sygnaturze & \ref{ch:services} \\
400 z \code{[ApiController]} & 422 niosące listę błędów pydantica &
  \ref{ch:services} \\
\code{IOptions<T>}, \code{appsettings.json} & \pkg{pydantic-settings} nad
  środowiskiem i \code{.env} & \ref{ch:services} \\
czasy życia DI & \code{Depends}, na żądanie; \code{yield} na zasięg;
  lifespan na singletony & \ref{ch:services} \\
middleware, pierwszy zarejestrowany jest zewnętrzny & ostatni
  zarejestrowany jest zewnętrzny & \ref{ch:services} \\
wątki Kestrela & workery uvicorna to procesy, równoważone na połączenie &
  \ref{ch:services} \\
\code{JsonPropertyName} & \code{serialization\_alias}; samo \code{alias}
  zmienia też wejście & \ref{ch:services} \\
Entity Framework & SQLAlchemy & \ref{ch:data} \\
\code{DbContext} & \code{Session} \dash{} i wygasza twoje obiekty przy
  commicie & \ref{ch:data} \\
\code{DbSet<T>}, LINQ & \code{select(T).where(...)}, które nie przyjmuje
  lambdy & \ref{ch:data} \\
\code{Include} & \code{selectinload}, przeciw temu samemu N+1 &
  \ref{ch:data} \\
\code{Add-Migration} & \code{alembic revision --autogenerate}, które
  musisz przeczytać & \ref{ch:data} \\
\code{Update-Database} & \code{alembic upgrade head} & \ref{ch:data} \\
DataTable & ramka \pkg{polars} & \ref{ch:data} \\
xUnit & pytest & \ref{ch:testing} \\
\code{[Fact]}, \code{[Theory]} & funkcja \code{test\_},
  \code{@pytest.mark.parametrize} & \ref{ch:testing} \\
\code{IClassFixture<T>} & funkcja fixture z zasięgiem &
  \ref{ch:testing} \\
Moq & \api{unittest.mock}, podmieniany tam, gdzie nazwa jest wyszukiwana &
  \ref{ch:testing} \\
\code{Assert.Equal} & zwykły \code{assert}, który pytest przepisuje, by
  sam się tłumaczył & \ref{ch:testing} \\
FluentAssertions & znowu zwykły \code{assert} & \ref{ch:testing} \\
\code{ILogger<T>} & \pkg{structlog} nad standardowym modułem
  \api{logging} & \ref{ch:observability} \\
\code{AsyncLocal<T>} & \api{contextvars} & \ref{ch:observability} \\
OpenTelemetry & OpenTelemetry i te same spany & \ref{ch:observability} \\
zamykanie \code{IHostedService} & lifespan uvicorna, który domyślnie czeka
  w nieskończoność & \ref{ch:observability} \\
\code{dotnet publish}, obraz runtime'owy & wieloetapowy Dockerfile i
  \code{uv sync --frozen --no-dev} & \ref{ch:observability} \\
\code{dotnet list package -{}-vulnerable} & \code{pip-audit}, nad
  \code{pylock.toml}, który eksportujesz & \ref{ch:observability} \\
\end{cheatsheet}

\section{Część V \dash{} Python do pracy z AI}
\label{sec:appA-part5}

\begin{cheatsheet}{C\#}{Python}{Rozdz.}
klient SDK & model pydantica, klient httpx i iterator strumieniowy &
  \ref{ch:aitoolkit} \\
\code{JsonSerializer.Deserialize} & \code{model\_validate\_json}, i jest
  szybsze niż wcześniejsze parsowanie & \ref{ch:aitoolkit} \\
ścisły deserializator & \code{strict=True}, które nic nie kosztuje &
  \ref{ch:aitoolkit} \\
\code{IAsyncEnumerable} po strumieniu & \code{async for} po strumieniu
  SDK & \ref{ch:aitoolkit} \\
projekt konsolowy na brudno & notatnik, a książka mówi, kiedy jest
  właściwy & \ref{ch:aitoolkit} \\
xUnit nad nagranym śladem & \code{pytest} nad \code{Trace}, z dwunastoma
  deterministycznymi asercjami & \ref{ch:traceassert} \\
\code{dotnet pack}, \code{nuget push} & \code{uv build}, \code{uv
  publish} z trusted publishing & \ref{ch:traceassert} \\
\end{cheatsheet}
